\documentclass{report}[10pt]
\begin{document}
\scriptsize
\begin{verbatim}
// Header file for Network, the abstract base class for neUROn2++ neural networks
//    which use some form of iterative training method.
//    All Networks have weights in some form which may be randomized,
//    a feed forward method, a training method, and a learning rate

#include "stdafx.h" // For MSVC, must be first!

#include "network.h"

// Default constructor
Network::Network()
{
   builtFlag = false; // the Network object hasn't yet been built
   weightsSetFlag = false; // weights not set yet

   // Good ideas for initial values
   randomLimit = 0.5; // a good initial random limit   
   biasFlag = true; // Networks should generally have bias nodes
   batchEpochFlag = true; // Networks should generally use batch/epoch training
   weightDecayFlag = true; // Networks should generally use weight decay
   decay = 5e-5; // lambda = 1/(2*sigma^2) where sigma = 100

   // Values that really should be set in the driver
   // The learning rate should be set in the driver to be 1 or 1/N depending on
   //    if training is off-line (batch/epoch) or on-line
   eta = 0.05;

   // Start with canonical backprop
   trainingType = 0;
}

// Default destructor
Network::~Network() { }

// Copy constructor
Network::Network( const Network& rhs )
{
   Network::copy( rhs ); // use the copy utility
}

// Overloaded = operator
Network& Network::operator= ( const Network& rhs )
{
   if ( &rhs != this ) // check for self-assignment
      Network::copy( rhs ); // use the copy utility
   
   return *this; // enables A = B = C
}

// Copy utility
void Network::copy( const Network& rhs )
{
   Iterative::copy( rhs ); // call immediate base object copy
   output = rhs.output;
   xs = rhs.xs;
   o = rhs.o;
   x = rhs.x;
   randomLimit = rhs.randomLimit;
   eta = rhs.eta;
   decay = rhs.decay;
   decayTerm = rhs.decayTerm;
   regularizer = rhs.regularizer;
   weightsSetFlag = rhs.weightsSetFlag;
   biasFlag = rhs.biasFlag;
   batchEpochFlag = rhs.batchEpochFlag;
   weightDecayFlag = rhs.weightDecayFlag;
   trainingType = rhs.trainingType;
   stackG = rhs.stackG;
   lastF = rhs.lastF;
   lastG = rhs.lastG;
}

// Set the learning rate
void Network::setEta( const double n )
{
   assert( n > 0 && n <= 1 );
   eta = n;
}

// Set the weight decay value
void Network::setDecay( const double n )
{
   assert( n >= 0 && n <= 1 );
   decay = n;
}

// Outputs to ostream reporting the accuracy of a Network Model
void Network::reportAccuracy( ostream& outputStream )
{
   unsigned r, // row iterator
      nInput = theData.getInput(); // number of inputs (easier on the eyes)
            
   double setError = 0, // set error for test & training sets
      accuracy = 0; // classification accuracy for multiple output sets
   
   if ( theData.getOutput() == 1 ) // accuracy report for 1-output Networks
   {
      if ( theData.trainLoaded() ) // final training set error
      {
         // Iterate through training exemplars
         for ( r = 0; r < Train.rows(); r++ )
         {
            // Forward propagate the training exemplar
            forward( Train, r );
            
            // Calculate error for single output and add to set error
            errorFunction E( theData.getTrainMatrix()( r, nInput ), o, x, errorType );
            setError += E.value();

            // If discrete, populate TwoSet object for accuracy report
            if ( theData.getDiscrete() && theData.getTrainTwoSet().loaded() )
               theData.getTrainTwoSet().test( r ) = o;
         }
         
         setError /= Train.rows(); // calculate training set error

         // Send resulting report to output stream
         outputStream << resetiosflags( ios::fixed )
            << setiosflags( ios::scientific ) << setprecision( 6 )
            << "Final " << errorLabel << " error in the training set = "
            << setError << "." << endl;
      }
      
      if ( theData.testLoaded() ) // final test set error
      {
         setError = 0; // reset set error
         
         // Iterate through test exemplars
         for ( r = 0; r < Test.rows(); r++ )
         {
            // Forward propagate the test exemplar
            forward( Test, r );
            
            // Calculate error for single output and add to set error
            errorFunction E( theData.getTestMatrix()( r, nInput ), o, x, errorType );
            setError += E.value();

            // If discrete, populate TwoSet object for accuracy report
            if ( theData.getDiscrete() && theData.getTestTwoSet().loaded() )
               theData.getTestTwoSet().test( r ) = o;
         }
         
         setError /= Test.rows(); // calculate test set error
         
         // Send resulting report to output stream
         outputStream << resetiosflags( ios::fixed )
            << setiosflags( ios::scientific ) << setprecision( 6 )
            << "Final " << errorLabel << " error in the test set = "
            << setError << "." << endl;
      }
      
      if ( theData.getDiscrete() ) // output TwoSet metrics report if discrete
         theData.metricsReport( outputStream );
   }   
   
   else // multiple outputs
   {
      if ( theData.trainLoaded() )
      {
         setError = 0; // reset set error

         // Iterate through training exemplars
         for ( r = 0; r < Train.rows(); r++ )
         {
            // Forward propagate the training exemplar
            forward( Train, r );
               
            // Calculate error for multiple output and add to set error
            errorFunction E( TrainOutput.row( r ), output, xs, errorType );
            setError += E.value();
         }            
            
         setError /= Train.rows(); // calculate training set error            
            
         // Send resulting report to output stream
         outputStream << resetiosflags( ios::fixed )
            << setiosflags( ios::scientific ) << setprecision( 6 )
            << "Final " << errorLabel << " error in the training set = "
            << setError << "." << endl;

         // If descrete, report final accuracy
         if ( theData.getDiscrete() )
         {
            accuracy = 0; // reset the accuracy

            for ( r = 0; r < Train.rows(); r++ ) // iterate through Train Matrix rows
            {
               forward( Train, r ); // calculate the network's output
               // Make output discrete
               func( output, threshold( theData.getThreshold() ), output );
               if ( output == TrainOutput.row( r ) ) // compare known to guess
                  accuracy++;
            }

            accuracy /= Train.rows(); // calculate the accuracy

            // Send formatted training classification accuracy to ostream
            outputStream << resetiosflags( ios::scientific )
               << setiosflags( ios::fixed ) << setprecision( 1 )
               << "Final classification accuracy in the training set = "
               << accuracy * 100 << "." << endl;
         }      
      }

      if ( theData.testLoaded() ) // for test set
      {
         setError = 0; // reset set error

         // Iterate through test exemplars
         for ( r = 0; r < Test.rows(); r++ )
         {
            // Forward propagate the test exemplar
            forward( Test, r );

            // Calculate error for multiple output and add to set error
            errorFunction E( TestOutput.row( r ), output, xs, errorType );
            setError += E.value();
         }            
            
         setError /= Test.rows(); // calculate test set error            
            
         // Send resulting report to output stream
         outputStream << resetiosflags( ios::fixed )
            << setiosflags( ios::scientific ) << setprecision( 6 )
            << "Final " << errorLabel << " error in the test set = "
            << setError << "." << endl;

         // If descrete, report final accuracy
         if ( theData.getDiscrete() )
         {
            accuracy = 0; // reset the accuracy

            for ( r = 0; r < Test.rows(); r++ ) // iterate through Test Matrix rows
            {
               forward( Test, r ); // calculate the network's output
               // Make output discrete
               func( output, threshold( theData.getThreshold() ), output );
               if ( output == TrainOutput.row( r ) ) // compare known to guess
                  accuracy++;
            }

            accuracy /= Test.rows(); // calculate the accuracy

            // Send formatted training classification accuracy to ostream
            outputStream << resetiosflags( ios::scientific )
               << setiosflags( ios::fixed ) << setprecision( 1 )
               << "Final classification accuracy in the test set = "
               << accuracy * 100 << "." << endl;
         }
      }
   }
}

// Outputs classification accuracies to ostream for Iterative table
//    Iterative.cpp has already checked that outputs are discrete
void Network::classAccuracy( ostream& outputStream )
{
   unsigned r; // row iterator
   
   if ( theData.getOutput() == 1 ) // for 1-output Networks
   {
      // Training TwoSet must exist
      if ( theData.trainLoaded() && theData.getTrainTwoSet().loaded() )
      {
         for ( r = 0; r < Train.rows(); r++ ) // iterate through Train Matrix rows
         {
            forward( Train, r ); // calculate the network's output
            // Set the DataSet training TwoSet member's test member to the output
            theData.getTrainTwoSet().test( r ) = o;
         }

         // Send formatted training classification accuracy to ostream
         outputStream << resetiosflags( ios::scientific )
            << setiosflags( ios::fixed );
         outputStream << "  " << setw( 10 ) << setfill( ' ' ) << setprecision( 1 )
            << theData.getTrainTwoSet().getClassAcc() * 100;
      }

      // Test TwoSet must exist
      if ( theData.testLoaded() && theData.getTestTwoSet().loaded() )
      {
         for ( r = 0; r < Test.rows(); r++ ) // iterate through Test Matrix rows
         {
            forward( Test, r ); // calculate the network's output
            // Set the DataSet test TwoSet member's test member to the output
            theData.getTestTwoSet().test( r ) = o;
         }

         // Send formatted training classification accuracy to ostream
         outputStream << resetiosflags( ios::scientific )
            << setiosflags( ios::fixed );
         outputStream << "  " << setw( 10 ) << setfill( ' ' ) << setprecision( 1 )
            << theData.getTestTwoSet().getClassAcc() * 100;
      }
   }
   
   else // multiple outputs
   {
      double accuracy; // the classification accuracy

      // For the training set
      if ( theData.trainLoaded() )
      {
         accuracy = 0; // reset the accuracy

         for ( r = 0; r < Train.rows(); r++ ) // iterate through Train Matrix rows
         {
            forward( Train, r ); // calculate the network's output
            // Make output discrete
            func( output, threshold( theData.getThreshold() ), output );
            if ( output == TrainOutput.row( r ) ) // compare known to guess
               accuracy++;
         }

         accuracy /= Train.rows(); // calculate the accuracy

         // Send formatted training classification accuracy to ostream
         outputStream << resetiosflags( ios::scientific )
            << setiosflags( ios::fixed );
         outputStream << "  " << setw( 10 ) << setfill( ' ' ) << setprecision( 1 )
            << accuracy * 100;
      }

      // For the test set
      if ( theData.testLoaded() )
      {
         accuracy = 0; // reset the accuracy

         for ( r = 0; r < Test.rows(); r++ ) // iterate through Test Matrix rows
         {
            forward( Test, r ); // calculate the network's output
            // Make output discrete
            func( output, threshold( theData.getThreshold() ), output );
            if ( output == TrainOutput.row( r ) ) // compare known to guess
               accuracy++;
         }

         accuracy /= Test.rows(); // calculate the accuracy

         // Send formatted test classification accuracy to ostream
         outputStream << resetiosflags( ios::scientific )
            << setiosflags( ios::fixed );
         outputStream << "  " << setw( 10 ) << setfill( ' ' ) << setprecision( 1 )
            << accuracy * 100;
      }
   }
}

// Utility to output Network specific parameters prior to an Iterative run
void Network::runHeader( ostream& outputStream )
{
   // A great place to initialize the actual weight multipliers for decay, as this needs
   //    only to be calculated once, right before TrainSet() is called iteratively
   regularizer = decay / 2; // the regularization term lambda
   decayTerm = 1 - ( eta * decay ); // 1 - 2*eta*lambda

   switch ( trainingType ) // the training algorithm
   {
      case 0:
         outputStream << "Training algorithm is canonical backpropagation" << endl;
         break;
      case 1:
         outputStream << "Training algorithm is conjugate gradient descent" << endl;
         break;
      case 2:
         outputStream << "Training algorithm is Shanno" << endl;
         break;
   }

   if ( batchEpochFlag ) // batch/epoch learning
      outputStream << "Training is off-line (batch/epoch is on)" << endl;
   else
      outputStream << "Training is on-line (batch/epoch is off)" << endl;

   outputStream << "Learning rate eta: " << eta << endl; // eta

   if ( weightDecayFlag ) // weight decay
      outputStream << "Weight decay term labmda: " << decay << endl;
   else
      outputStream << "Weight decay is off" << endl;
}

// Returns maximum gradient of this Network object
double Network::getGradMax()
{
   // If stackG is not created in trainSet() (e.g. canonical backprop),
   //    then one must be created.
   // NOTE: this implies that all other types of Networks create stackG.
   //    If you are making a Network that does not create stackG, you must
   //    change this if statement so that it will be packed.
   if ( trainingType == 0 )
      pack();

   return maxabs( stackG );
}

// The master training engine, takes training type and iteration as arguments,
//    calls the appropriate training algorithm
void Network::engine( unsigned type, unsigned t )
{
   switch ( type ) // choose training algorithm
   {
      case 1:
         CGD( t ); // conjugate gradient descent
         break;
      case 2:
         shanno( t ); // Shanno algorithm
         break;
   }
}

// Congugate gradient descent, Golden pp. 221-222
void Network::CGD( unsigned t )
{
   pack(); // creates gradient vector stackG from gradient structure

   if ( ( t == 0 ) || ( t == df() ) ) // Golden, step 1
\end{verbatim}
\hspace{24 pt}//\hspace{4 pt}$\vec{f}(t)=-\vec{g}_t$
\begin{verbatim}
      // Note that lastF refers to *negative* F in Golden's book,
      // So for the purposes here, return stackG,
\end{verbatim}
\hspace{24 pt}//\hspace{4 pt}and remember $\vec{g}_{t-1}$ for next iteration
\begin{verbatim}
      lastG = stackG;

   else // Golden, step 2
   {
\end{verbatim}
\hspace{24 pt}//\hspace{4 pt}$\vec{u}_t=\vec{g}_t-\vec{g}_{t-1}$
\begin{verbatim}
      vector< double > u = stackG - lastG;
\end{verbatim}
\hspace{24 pt}//\hspace{4 pt}Remember $\vec{g}_{t-1}$ for next iteration
\begin{verbatim}
      lastG = stackG;

      // Since lastF refers to *negative* F in Golden's book
      // D will thus have reversed sign (-) from Golden's book
      double D = dotprod( u, lastF );

      if ( D != 0 ) // catch division by zero
      {
\end{verbatim}
\hspace{36 pt}//\hspace{4 pt}$b_t=\frac{\vec{u}^T\vec{g_t}}{\vec{u}^T_t\vec{f}(t-1)}$,
$b_t$ will thus have reversed sign (-) from Golden's book
\begin{verbatim}
         double b = dotprod( u, stackG ) / D;
\end{verbatim}
\hspace{36 pt}//\hspace{4 pt}$\vec{f}(t)=-\vec{g}_t+b_t\vec{f}(t-1)$
\begin{verbatim}
         // with the sign notes from above
         stackG -= ( lastF * b );
      }
   }

\end{verbatim}
\hspace{12 pt}//\hspace{4 pt}stackG was loaded with $\vec{f}$
\begin{verbatim}
   lastF = stackG;

   unpack(); // return gradient structure from stackG
}

// Shanno algorithm, Golden pp. 217-218
void Network::shanno( unsigned t )
{
   pack(); // creates gradient vector stackG from gradient structure

   if ( ( t == 0 ) || ( t == df() ) ) // Golden, step 1
\end{verbatim}
\hspace{24 pt}//\hspace{4 pt}$\vec{f}(t)=-\vec{g}_t$
\begin{verbatim}
      // Note that lastF refers to *negative* F in Golden's book,
      // So for the purposes here, return stackG,
\end{verbatim}
\hspace{24 pt}//\hspace{4 pt}and remember $\vec{g}_{t-1}$ for next iteration
\begin{verbatim}
      lastG = stackG;

   else // Golden, step 2
   {
\end{verbatim}
\hspace{24 pt}//\hspace{4 pt}$\vec{u}_t=\vec{g}_t-\vec{g}_{t-1}$
\begin{verbatim}
      vector< double > u = stackG - lastG;
\end{verbatim}
\hspace{24 pt}//\hspace{4 pt}Remember $\vec{g}_{t-1}$ for next iteration
\begin{verbatim}
      lastG = stackG;

      // Since lastF refers to *negative* F in Golden's book
      // D will thus have reversed sign (-) from Golden's book
      double D = dotprod( lastF, u ); // denominator for all coefficients

      if ( D != 0 ) // catch division by zero
      {
\end{verbatim}
\hspace{36 pt}//\hspace{4 pt}$a_t=\frac{\vec{f}(t-1)^T\vec{g_t}}{\vec{f}(t-1)^T\vec{u}_t}$,
$a_t$ will thus have same sign as Golden's book
\begin{verbatim}
         double a = dotprod( lastF, stackG ) / D;
\end{verbatim}
\hspace{36 pt}//\hspace{4 pt}$b_t=\frac{\vec{u}^T\vec{g_t}}{\vec{f}(t-1)^T\vec{u}_t}$,
$b_t$ will thus have reversed sign (-) from Golden's book
\begin{verbatim}
         double b = dotprod( u, stackG ) / D;
\end{verbatim}
\hspace{36 pt}//\hspace{4 pt}$c_t=\gamma+\frac{|\vec{u}_t|^2}{\vec{f}(t-1)^T\vec{u}_t}$,
2nd term in $c_t$ will have reversed sign (-) from Golden's book
\begin{verbatim}
         double c = eta - ( squared( u ) / D );

\end{verbatim}
\hspace{36 pt}//\hspace{4 pt}$\vec{f}(t)=-\vec{g}_t+a_t\vec{u}_t+(b_t-c_ta_t)\vec{f}(t-1)$
\begin{verbatim}
         // with the sign notes from above
         stackG -= ( u * a );
         lastF *= ( b  + ( c * a ) );
         stackG -= lastF;
      }
   }

\end{verbatim}
\hspace{12 pt}//\hspace{4 pt}stackG was loaded with $\vec{f}$
\begin{verbatim}
   lastF = stackG;

   unpack(); // return gradient structure from stackG
}
\end{verbatim}
\end{document}
